\section{\ref{sec:chemoutput}장. 화학적 출력}

심장으로 가는 두 가닥의 전선과 그 속도 차이, 호르몬 캐스케이드의 음의 되먹임, 분량 빈도에 의한 부호, 일주기 발진기와 빛에 의한 맞춤은 직접 측정되었다. $K<8$이라는 경계는 본문에서 유도한 제어 이론의 정리다. 캐스케이드의 되먹임 자체가 맥동을 만든다는 것은 이를 강하게 뒷받침하는 실험이 있는 가설이다.

{\footnotesize
\setlength{\tabcolsep}{3pt}\renewcommand{\arraystretch}{1.15}
\begin{longtable}{>{\raggedright}p{0.5cm}>{\raggedright}p{4.0cm}>{\raggedright}p{5.6cm}>{\raggedright}p{2.3cm}>{\raggedright\arraybackslash}p{1.7cm}}
\toprule
No. & 명제 & 실험과 측정한 것 & 출처 & 상태 \\
\midrule\endfirsthead
\toprule
No. & 명제 & 실험과 측정한 것 & 출처 & 상태 \\
\midrule\endhead
1 & 심장 페이스메이커는 스스로 1분에 약 $100$번 뛴다. 안정 시에는 브레이크가 밟혀 있어 박동수가 보통 $60$--$70$ 정도다 & 사람, 심장으로 가는 두 전선을 모두 차단: 고유 박동수는 안정 시 박동수보다 높고 나이가 들수록 줄며, 성인에서 분당 약 $100$회다. 따라서 안정 시에는 브레이크가 우세하다. 안정 시 박동수 $60$--$70$은 전형적인 값으로 따로 인용하지 않았다. & \cite{JoseCollison1970ihr} & 측정됨 \\
2 & 가속 페달 \nt{NORA}와 브레이크 \nt{ACET}은 저역 통과 필터이며 브레이크가 더 빠르다~\eqref{eq:gasbrake} & 개, 심장의 미주 신경 또는 교감 신경을 주파수 변조 자극하고 심방 박동수까지의 전달 함수를 측정: 두 경로 모두 저역 통과 필터이고, 교감 경로는 차단 주파수가 훨씬 낮으며 약 $1.7\,\text{s}$의 순수 지연이 더해진다. 특성은 평균 긴장도에 따라 달라지므로 선형식~\eqref{eq:gasbrake}은 근사다. & \cite{Berger1989} & 측정됨 \\
3 & 브레이크가 빠른 것은 \nt{ACET}이 2차 전달자 없이 칼륨 채널을 열기 때문이고, \nt{NORA}는 연쇄 반응을 거친다 & 심방 세포의 막 조각: 아세틸콜린이 여는 칼륨 채널은 수용체와 결합한 G 단백질의 $\beta\gamma$ 소단위가 세포 안 2차 전달자 없이 연다. cAMP를 거치는 \nt{NORA} 경로는 여기서 인용하지 않았다. & \cite{Logothetis1987gbg} & 측정됨 \\
4 & 혈액은 약 1분 만에 몸을 한 바퀴 돈다. 혈류 속 \nt{ADRE}는 수용체가 있는 모든 기관에 닿는다 & 일반적인 크기 수준의 추정. 본문에서도 그렇게 서술했고 출처는 인용하지 않았다. & --- & 추정 \\
5 & 호르몬 캐스케이드는 음의 되먹임으로 감싸여 있다. 최종 호르몬이 앞 단계를 억제한다 & 실험 리뷰: 부신 피질 호르몬은 뇌하수체의 ACTH 분비와 시상하부의 방출 호르몬 분비를 빠른, 중간, 느린 세 시간 범위에서 억제한다. & \cite{KellerWood1984fb} & 측정됨 \\
6 & 같은 단 세 개는 $K<8$에서 안정하고~\eqref{eq:cascadestable}, 경계에서의 주기는 $2\pi\tau/\sqrt3\approx3.6\tau$ & 본문에서 유도: 주파수 $\sqrt3/\tau$에서 세 단은 위상을 $180^\circ$ 늦추고 진폭을 $8$분의 1로 줄인다. & 본문에서 유도 & 이론 \\
7 & 농도 오차는 $1/(1+K)$이므로 이런 조절기는 약 $10\,\%$보다 정확하게 유지하지 못한다(명제~\ref{prop:cascade}) & 비례 되먹임에 대한 6행의 따름정리. 실제 축에는 여러 단계와 여러 시간 범위에 되먹임이 있으므로(5행), 이 경계는 모델~\eqref{eq:cascade}에 대한 것이며 측정된 것이 아니다. & 본문에서 유도 & 이론 \\
8 & $K=4$와 $7$에서 농도가 자리 잡고, $K=12$에서는 주기 $3.7$--$3.9\,\tau$의 일정한 맥동(그림~\ref{fig:cascade}) & \texttt{models/\allowbreak chem\_\allowbreak models.py}로 계산: $K=12$에서 연속된 주기 $3.9$, $3.7$, $3.8$, $3.7\,\tau$; $K=4$에서 계산 끝의 진폭 $0.003$. & \texttt{models/\allowbreak chem\_\allowbreak models.py} & 모델 \\
9 & 스트레스 호르몬은 약 1시간에 한 번 맥동으로 분비된다. 맥동은 별도의 발진기 없이 되먹임 자체가 만든다 & 쥐, 방출 호르몬을 일정하게 주입: ACTH와 코르티코스테론이 생리적인, 거의 1시간 주기로 오르내린다. 맥동에 시상하부의 리듬 입력은 필요 없다. 지연된 되먹임이 있는 뇌하수체--부신 모델이 이런 진동을 낸다. & \cite{Walker2012glc, Walker2010} & 가설 \\
10 & 수신자는 농도가 아니라 분량의 빈도를 읽는다 & 생식샘 자극 호르몬 방출 호르몬을 스스로 분비하지 못하는 원숭이: 일정한 주입은 뇌하수체 호르몬 분비를 회복시키지 못하고, 1시간에 한 번 분량으로 주면 회복된다. 다시 일정한 주입으로 바꾸면 반응이 다시 사라진다. 고역 통과 필터로서의 수용체 피로는 해석이다. & \cite{Belchetz1978} & 측정됨 \\
11 & 분량의 빈도가 다르면 같은 샘의 서로 다른 세포에 다르게 작용한다 & 같은 원숭이: 분량을 1시간에 $2$--$5$회로 늘리면 두 뇌하수체 호르몬이 모두 줄고, $3$시간에 1회로 줄이면 하나(LH)는 줄고 다른 하나(FSH)는 늘어난다. 따라서 둘의 비는 빈도가 정한다. & \cite{Wildt1981gnrh} & 측정됨 \\
12 & 일주기 리듬은 몇 시간 지연되는 세포 안의 음의 되먹임이 만든다 & 리뷰: 시계 유전자의 단백질이 지연을 두고 자기 유전자의 전사를 억제한다. 전사와 번역으로 이루어진 루프가 약 하루의 주기를 낸다. & \cite{TakahashiJS2017} & 측정됨 \\
13 & 주 발진기는 수만 개의 뉴런으로 이루어진 무리이며 몸의 나머지를 동기화한다 & 리뷰: 시교차상핵이 주 일주기 발진기다. 배양한 개별 뉴런은 독립된 발진기이고, 네트워크가 이들을 동기화한다. 생쥐에서는 한쪽에 약 $10^4$개의 뉴런이 있다. & \cite{Welsh2010scn} & 측정됨 \\
14 & 사람의 고유 주기는 $24.18$시간 & 하루와 분리된 일정 아래 어두운 빛에서 지낸 사람: 멜라토닌, 체온, 코르티솔 리듬의 주기는 젊은 사람과 노인 모두 평균 $24.18$시간이고 편차가 좁다. & \cite{Czeisler1999} & 측정됨 \\
15 & 빛은 위상 고정 루프처럼 발진기를 맞춘다~\eqref{eq:pll}. 하루 약 $11$~min의 이동이 필요하다 & 사람, 하루 중 여러 시각에 $6.7$시간의 밝은 빛을 한 번 비춤: 진폭 $5$시간의 1형 위상 반응 곡선. 체온 최저점 이전에는 지연, 이후에는 앞당김. 식~\eqref{eq:pll}은 표준 모델이고, $11$~min $=0.18$시간은 산술이다. & \cite{Khalsa2003prc} & 측정됨 \\
16 & 빛이 없으면 맞춤이 없고 리듬은 고유 주기로 흘러간다 & 빛을 전혀 지각하지 못하고 일상 사회에서 사는 전맹인: 약 절반에서 멜라토닌과 코르티솔 리듬이 하루와 맞지 않는 자기 주기로 간다. & \cite{Sack1992blind} & 측정됨 \\
\bottomrule
\end{longtable}}
